\documentclass[11pt]{article}
\usepackage[a4paper,margin=27mm]{geometry}
\usepackage{amsmath,amssymb,amsthm,hyperref}
\newtheorem{theorem}{Theorem}
\newtheorem{lemma}[theorem]{Lemma}
\newcommand{\Lie}{\operatorname{Lie}}
\newcommand{\R}{\mathbb R}
\title{A factorial-scale barrier even for error-dependent relabeling averages}
\author{Mathematical proof reviewed independently; not formally verified}
\date{30 September 2026}
\begin{document}
\maketitle
Let $V=\Lie_n$ be the real multilinear degree-$n$ free Lie
representation on $n$ distinct generators. It has dimension $(n-1)!$.
For $E\in V$ let $q(E)$ be the least support size of a probability
measure $\mu$ on $S_n$ such that
$\sum_g\mu(g)gE=0$. Such measures always exist, by uniform averaging.
Unlike a universal annihilator, the measure here may depend on $E$.

\begin{theorem}\label{main}
For $n\ge7$ put
\[
 Q_n=\left\lfloor\sqrt{\frac{n!}{2n^5}}\right\rfloor.
\]
If $Q_n\ge1$, the set $\{E\in V:q(E)\le Q_n\}$ is closed and has
Lebesgue measure zero. Its complement contains an element with integral
coordinates in the standard fixed-last-letter bracket basis.

In particular, for every sufficiently large $n$, there is a positive
word $w$ with equal letter counts, fair on every subset of at most
$n-1$ letters, such that cancelling its top-degree defect by positive
averaging of relabeled copies requires more than $Q_n$ distinct
relabelings. Thus the worst-case support requirement is
\[
 \exp\bigl((1/2+o(1))n\log n\bigr).
\]
\end{theorem}

\begin{lemma}[Rank of a positive relabeling average]\label{rank}
For $n\ge7$, if $\mu$ is supported on at most $q$ permutations, then
\[
 \operatorname{rank}(T_\mu|_V)\ge\frac{n!/q-2}{n^5},
 \qquad T_\mu=\sum_g\mu(g)g.
\]
\end{lemma}
\begin{proof}
On the regular representation of $S_n$, every group element is
orthogonal, so $\|T_\mu\|_{\mathrm{op}}\le1$. Its squared Frobenius
norm is $n!\sum_g\mu(g)^2\ge n!/q$. Since each singular value is at
most one, its regular-representation rank is at least $n!/q$.

For an irreducible complex representation $S^\lambda$, let $d_\lambda$
be its dimension and $m_\lambda$ its multiplicity in $\Lie_n$.
Klyachko's theorem gives $m_\lambda\ge1$ for every shape except trivial
and sign when $n\ge7$. Moreover,
\[
 m_\lambda\ge d_\lambda/n^5
 \quad\text{for every such shape}.\tag{1}
\]
Indeed, this is immediate when $d_\lambda\le n^5$; when
$d_\lambda>n^5$, Swanson \cite[Theorem 1.9]{swanson}, together
with the Lie multiplicity interpretation, gives
$m_\lambda/d_\lambda>1/n-1/n^2\ge1/n^5$.

If $r_\lambda$ is the rank of $\sum_g\mu(g)\rho_\lambda(g)$,
complexification and semisimplicity give
\[
 \operatorname{rank}_{\mathrm{reg}}T_\mu
   =\sum_\lambda d_\lambda r_\lambda,
 \qquad
 \operatorname{rank}_{V}T_\mu
   =\sum_\lambda m_\lambda r_\lambda.
\]
The two omitted one-dimensional representations contribute at most two
to the first sum. Inequality (1) proves the result.
\end{proof}

\begin{lemma}[A parameterized-kernel observation]\label{kernels}
Let $A(t)$ be a linear operator on $\R^d$ whose entries depend linearly
on $t$ in a compact simplex of dimension at most $q-1$. Suppose
$\operatorname{rank}A(t)\ge q$ throughout the simplex. Then
\[
 \mathcal B=\{x:\text{some }t\text{ satisfies }A(t)x=0\}
\]
is closed and has Lebesgue measure zero.
\end{lemma}
\begin{proof}
Closedness follows by taking a convergent subsequence of parameters
in the compact simplex whenever $x_j\to x$ and $A(t_j)x_j=0$.
For the measure statement, cover the parameter simplex by the finitely
many sets where a specified $q\times q$ minor is nonzero. On each such
set, $q$ of the coordinates of a kernel vector are rational, hence smooth,
functions of the parameters and the remaining $d-q$ coordinates.
Ignoring any remaining equations only enlarges the resulting set.
The total number of free real parameters is at most
$(q-1)+(d-q)=d-1$. Restricting to countably many bounded pieces away
from a zero denominator makes each parameterization Lipschitz, and a
Lipschitz image of a subset of $\R^{d-1}$ has $d$-dimensional measure
zero. The finite union of charts has the same property.
\end{proof}

\begin{proof}[Proof of Theorem~\ref{main}]
Put $q=Q_n\ge1$. Then $n!\ge2n^5q^2\ge n^5q^2+2q$, so
Lemma~\ref{rank} gives $\operatorname{rank}(T_\mu|_V)\ge q$ for every
measure supported on at most $q$ permutations.
For each fixed support of size $q$, its probability measures form a
compact $(q-1)$-simplex. Lemma~\ref{kernels} shows that the errors it
can cancel form a closed null set. There are finitely many supports;
smaller supports are included by allowing zero weights. Their union is
therefore closed and null. Its nonempty open complement contains a
rational point in the rational bracket basis. Multiplication by a
common denominator preserves membership in this complement, giving
an integral error $E$ with $q(E)>Q_n$.

We now realize this error by an actual positive word. In the reduced
nilpotent group, every degree-$n$ bracket is central and its group
commutator is exactly its exponential. Thus an integral linear
combination $E$ of the fixed-last-letter brackets gives a signed word
with signature $\exp(E)=1+E$.
The positive fraction-clearing argument in
\path{../positive_saturation/eventual_primorial_saturation.tex}
then supplies an equal-count fair signature $T_L=\exp(LH)$, where
$H=X_1+\cdots+X_n$, such that $\exp(E)T_L$ is represented by a
positive word $w$. Explicitly, write the signed element as $uv^{-1}$
with positive $u,v$; complete $v$ to a fair positive word $vc=T_L$,
and use $uc$. A further fair factor can make $L>0$ if necessary.
Centrality gives
\[
 \log S(w)=LH+E.
\]
Hence all restrictions to fewer than $n$ letters are fair and every
letter occurs exactly $L$ times.

For relabeled copies of $w$, the degree-$n$ defects are $gE$.
Their positive weighted average vanishes exactly when its normalized
weights define a measure annihilating $E$, which needs more than $Q_n$
relabelings. If integer copies and common integer blowups are used,
a blowup by $a$ multiplies the defect by $a^n>0$; these still produce
positive weights and cannot circumvent the support bound.
\end{proof}

\paragraph{Scope.}
This is a worst-case obstruction for cancelling the final defect of an
arbitrary almost-fair seed by relabeling. It does not show that every
seed needs many relabelings, does not rule out exploiting particular
seed structure, and gives no lower bound on the optimal size of a fair
family constructed by other means. The representation estimates are
classical; historical novelty of this application has not been audited.

\begin{thebibliography}{9}
\bibitem{swanson} Joshua P.~Swanson, \emph{On the existence of tableaux
with given modular major index}, Algebraic Combinatorics 1 (2018), 3--21,
\href{https://doi.org/10.5802/alco.4}{doi:10.5802/alco.4}.
\end{thebibliography}
\end{document}
